\section{Proposed AC-MOT Framework}\label{sec:method}
This section presents the proposed AC-MOT framework for adaptive multi-object tracking in aerial videos. The framework dynamically adjusts the detector operating point according to changing scene conditions without modifying or retraining the detector or tracker. It combines lightweight scene analysis, scene complexity estimation, and a three-state control mechanism to select suitable detector configurations during online tracking. The following subsections describe the operating point, the scene cues, the scene complexity index, the three-state controller, the action map, and how each value of the controller was chosen. Table~\ref{tab:symbols} lists the symbols used in this section.

\subsection{Overview and operating point}
Figure~\ref{fig:pipeline} shows where AC-MOT sits in the pipeline. For each
frame, scene analysis and the AC-MOT controller run first and choose the
detector settings. The detector and the tracker then run as usual with these
settings. The detector settings at frame $t$ are called the operating point
\begin{equation}
\theta_t = (c_t, \eta_t, r_t),
\label{eq:theta}
\end{equation}
which has three parts:
\begin{itemize}
\item $r_t$ is the input resolution, the long side of the resized image in
      pixels. A larger input shows small objects better but takes more time.
\item $\eta_t$ is the NMS IoU threshold. When two boxes overlap with an
      intersection over union (IoU) above $\eta_t$, the box with the lower
      score is removed. A lower $\eta_t$ removes more overlapping boxes.
\item $c_t$ is the score threshold. Boxes with a detector score below $c_t$
      are removed before they reach the tracker. A higher $c_t$ passes fewer
      but more reliable boxes.
\end{itemize}
The detector weights and the tracker are frozen; AC-MOT changes only
$\theta_t$.

The controller uses only information that is available at run time:
\begin{equation}
\theta_t = \pi\left(I_t, \mathcal{B}_{1:t-1}\right),
\label{eq:policy}
\end{equation}
where $\pi$ is the control rule described in this section, $I_t$ is the image
of the current frame, and $\mathcal{B}_{1:t-1}$ are the boxes that the tracker
reported for the earlier frames $1$ to $t-1$. In words, the decision for frame
$t$ uses the current image and the past tracker output. The tracker output of
frame $t$ is used only from frame $t+1$ on. No ground truth, no future frame,
and no detector score enters the decision.

\begin{figure*}[t]
\centering
\begin{tikzpicture}[
  font=\footnotesize,
  node distance=4mm,
  stage/.style={draw, rounded corners=2pt, align=center, text width=20mm,
                minimum height=11mm, inner sep=2pt},
  ctrl/.style={stage, fill=orange!10},
  host/.style={stage, fill=blue!7},
  note/.style={align=center, text width=22.5mm, font=\footnotesize, inner sep=1pt},
  arr/.style={-{Stealth[length=2mm]}, thick},
  fb/.style={-{Stealth[length=2mm]}, thick, dashed, green!40!black}]
\node[stage] (frame) {Frame $I_t$};
\node[ctrl, right=of frame] (scene) {Scene analysis};
\node[ctrl, right=of scene] (ctrl) {AC-MOT\\controller};
\node[host, right=of ctrl] (det) {Detector\\(frozen)};
\node[host, right=of det] (trk) {Tracker\\(frozen)};
\node[stage, right=of trk] (out) {Tracks of\\frame $t$};
\draw[arr] (frame) -- (scene);
\draw[arr] (scene) -- (ctrl);
\draw[arr] (ctrl) -- node[above, font=\footnotesize] {$\theta_t$} (det);
\draw[arr] (det) -- (trk);
\draw[arr] (trk) -- (out);
\node[note, below=2mm of scene] {edges, brightness, blur of $I_t$; crowd and tiny ratio from tracks of $t-1$};
\node[note, below=2mm of ctrl] {SCI, scene label, recovery, hysteresis, dwell, state};
\node[note, below=2mm of det] {resolution $r_t$, NMS IoU $\eta_t$, score threshold $c_t$};
\node[note, below=2mm of trk] {ByteTrack or OATrack, unchanged};
\draw[fb] (trk.north) -- ++(0,5mm) -| node[pos=0.25, above, font=\footnotesize]
  {tracked boxes, used from frame $t+1$} (scene.north);
\end{tikzpicture}

\caption{AC-MOT in a tracking-by-detection pipeline. Scene analysis and the
controller choose the operating point $\theta_t=(c_t,\eta_t,r_t)$ before the
detector runs. The frozen detector and tracker are not modified. Tracked boxes
of frame $t$ feed the scene analysis of later frames (dashed arrow).}
\label{fig:pipeline}
\end{figure*}

\begin{table}[t]
\caption{Symbols computed at run time. Fixed values are listed in
Table~\ref{tab:constants}.}
\label{tab:symbols}
\centering\footnotesize
\setlength{\tabcolsep}{4pt}
\begin{tabular}{@{}l>{\raggedright\arraybackslash}p{55mm}@{}}
\toprule
Symbol & Meaning\\
\midrule
$t$ & frame index\\
$j$ & analysis step (frames $1, 11, 21, \dots$)\\
$I_t$ & image of frame $t$\\
$r_t$, $\eta_t$, $c_t$ & input resolution, NMS IoU threshold, score threshold
  used at frame $t$, Eq.~\eqref{eq:theta}\\
$\mathcal{B}_{1:t-1}$ & boxes reported by the tracker for frames $1$ to $t-1$\\
$e_t$, $\beta_t$, $\nu_t$ & edge density, brightness, and sharpness of $I_t$\\
$N_j$ & number of tracked boxes in frame $t-1$\\
$a_i$ & area of tracked box $i$ in pixels\\
$C_j$ & crowd cue, Eq.~\eqref{eq:crowd}\\
$T_j$ & tiny-object ratio, Eq.~\eqref{eq:tiny}\\
$S_j$ & scene complexity index (SCI), Eq.~\eqref{eq:sci}\\
$H_j$ & the last seven analysis steps\\
$\ell_j$ & scene label, Eq.~\eqref{eq:label}\\
$g_j$ & target state, Eqs.~\eqref{eq:target}--\eqref{eq:hold}\\
$q_j$ & controller state: LOW, MEDIUM, or HIGH, Eq.~\eqref{eq:dwell}\\
$P_j$, $d_j$ & peak object count and drop counter,
  Eqs.~\eqref{eq:drop}--\eqref{eq:peak}\\
$t_c$ & frame of the last state change\\
\bottomrule
\end{tabular}
\end{table}

\subsection{Scene cues}
Scene analysis does not run on every frame. It runs on frame 1 and then once
every 10 frames ($t=1$ or $t \bmod 10 = 1$, that is, frames 1, 11, 21,
\dots). Each of these frames is an analysis step $j$. Between two analysis
steps the controller keeps its last decision. The stride of 10 frames was
selected by a temporal sweep on validation data (Sect.~\ref{sec:constants}).
Two kinds of cues are computed at each step.

\emph{Image cues of the current frame.} The frame is reduced to a quarter of
its width and height and converted to gray levels, so that the cues are cheap
to compute. Three cues are then measured:
\begin{itemize}
\item the edge density $e_t$, the fraction of pixels marked as edges by the
      Canny detector (lower and upper thresholds 50 and 120); many edges
      indicate a cluttered scene;
\item the brightness $\beta_t$, the mean gray level from 0 (black) to 255
      (white); a low value indicates a dark or night scene;
\item the sharpness $\nu_t$, the variance of the Laplacian of the gray image;
      a low value indicates a blurred frame.
\end{itemize}

\emph{Box cues of the previous frame.} Let $N_j$ be the number of boxes that
the tracker reported for frame $t-1$, and let $a_i$ be the area of box $i$ in
pixels of the original image. The crowd cue and the tiny-object ratio are
\begin{align}
C_j&=\min\!\left(\frac{N_j}{30},1\right),\label{eq:crowd}\\
T_j&=\frac{1}{N_j}\sum_{i=1}^{N_j}\mathbf{1}\!\left[a_i<32^2\right].\label{eq:tiny}
\end{align}
$C_j$ is the number of tracked objects divided by 30 and capped at 1, so every
frame with 30 or more objects has $C_j=1$. $T_j$ is the fraction of tracked
boxes smaller than $32\times32$ pixels, which is the size limit of small
objects in the COCO benchmark \cite{lin2014coco}. The indicator
$\mathbf{1}[\cdot]$ is 1 when its condition holds and 0 otherwise, and $T_j=0$
when $N_j=0$.

\subsection{Scene complexity index and scene label}
The scene complexity index (SCI) measures how crowded the recent frames are.
It is the mean crowd cue of the last seven analysis steps:
\begin{equation}
S_j=\frac{1}{|H_j|}\sum_{k\in H_j} C_k,
\label{eq:sci}
\end{equation}
where $H_j$ holds the steps $j-6,\dots,j$ (fewer at the start of a video) and
$|H_j|$ is their number. With a stride of 10 frames, the seven steps cover
about the last 60 frames, so a single unusual frame changes $S_j$ only a little. The
window of seven steps was selected by the same temporal sweep as the stride
(Sect.~\ref{sec:constants}). Because $C_j\le1$, $S_j$ lies between 0 and 1.
While fewer than 30 objects are tracked, $S_j$ is simply the recent average
number of tracked objects divided by 30.

The SCI uses the crowd cue alone. During development, all 31 combinations of
the five cues (crowd, tiny ratio, edges, brightness, and blur) were tested as
SCI inputs on the calibration sequences, and the crowd cue alone was selected
(Sect.~\ref{sec:modern}). The other four cues are still used as checks on the
state through the scene label
\begin{equation}
\ell_j=\begin{cases}
\text{night}   & \text{if } \beta_t<80,\\
\text{blur}    & \text{else if } \nu_t<180,\\
\text{tiny}    & \text{else if } T_j>0.50,\\
\text{crowded} & \text{else if } C_j>0.65 \text{ or } e_t>0.13,\\
\text{clear}   & \text{otherwise.}
\end{cases}
\label{eq:label}
\end{equation}
The conditions are checked from top to bottom, and the first one that holds
gives the label. For example, $T_j>0.50$ means that more than half of the
tracked boxes are tiny, and $C_j>0.65$ means that at least 20 objects are
tracked.

\subsection{Three-state controller}
The controller is always in one of three states: LOW, MEDIUM, or HIGH. The
names describe how difficult the scene is. They do not rank the input size:
in the final action map the HIGH state uses the smallest input
(Table~\ref{tab:actions}). The state changes only at analysis steps, in four
steps (Algorithm~\ref{alg:controller}).

\emph{Step 1: target state.} The SCI and the scene label give a target state
\begin{equation}
g_j=\begin{cases}
\text{HIGH}   & \text{if } S_j>\tau_h \text{ or } T_j>0.50,\\
\text{MEDIUM} & \text{else if } S_j>\tau_m \text{ or}\\
              & \quad \ell_j\in\{\text{crowded},\text{tiny}\},\\
\text{LOW}    & \text{otherwise,}
\end{cases}
\label{eq:target}
\end{equation}
where $\tau_m=0.297$ and $\tau_h=0.511$ are the two SCI thresholds (exact
values in Table~\ref{tab:constants}). They were selected by the joint search
on the calibration sequences (Sect.~\ref{sec:constants}). Since $S_j$ is the
recent average number of tracked objects divided by 30, the SCI conditions
give MEDIUM above about 9 objects ($30\tau_m=8.9$) and HIGH above about 15
objects ($30\tau_h=15.3$).

\emph{Step 2: recovery probe.} When the number of tracked objects suddenly
collapses, the controller tries the HIGH action for a short time. It keeps a
slowly decaying peak $P_j$ of the object count and a drop counter $d_j$:
\begin{align}
d_j&=\begin{cases}
d_{j-1}+1 & \text{if } P_{j-1}\ge5 \text{ and}\\
          & \quad N_j<0.40\,P_{j-1},\\
0         & \text{otherwise,}
\end{cases}\label{eq:drop}\\
P_j&=\max\left(N_j,\lfloor 0.8\,P_{j-1}\rfloor\right),\label{eq:peak}
\end{align}
with $P_0=d_0=0$. A collapse is detected when at least 5 objects were seen
before and fewer than 40\% of them remain. The peak loses 20\% at each step,
so an old peak is forgotten after a few steps. While $1\le d_j\le3$, the
target $g_j$ is set to HIGH, so the probe lasts at most three analysis steps
(30 frames).

\emph{Step 3: hysteresis.} A move to a lower state is held while the scene is
still close to difficult:
\begin{equation}
g_j\leftarrow\begin{cases}
\text{HIGH}   & \text{if } q_{j-1}=\text{HIGH} \text{ and}\\
              & \quad (S_j>0.50 \text{ or } T_j>0.40),\\
\text{MEDIUM} & \text{if } q_{j-1}=\text{MEDIUM},\\
              & \quad g_j=\text{LOW} \text{ and } S_j>0.25,\\
g_j           & \text{otherwise.}
\end{cases}
\label{eq:hold}
\end{equation}
The hold values lie just below the values that lead into each state
($0.50<\tau_h$, 0.40 below the tiny guard of 0.50, and $0.25<\tau_m$),
so the state does not jump back and forth when the scene stays near a
threshold. Upward moves are not held.

\emph{Step 4: dwell.} After a change, the state is kept for at least 30 frames,
that is, three analysis steps:
\begin{equation}
q_j=\begin{cases}
g_j     & \text{if } t-t_c\ge30,\\
q_{j-1} & \text{otherwise,}
\end{cases}
\label{eq:dwell}
\end{equation}
where $t_c$ is the frame of the last state change. The controller starts in
LOW with $t_c=1$, so the first change is possible at frame 31. Whenever
$q_j\ne q_{j-1}$, $t_c$ is set to $t$.

The state $q_j$ then selects the operating point from the action map, and the
detector uses it until the next analysis step.

\begin{algorithm}[t]
\caption{AC-MOT controller for one video}\label{alg:controller}
\begin{algorithmic}[1]
\State $q\gets$ LOW; $H\gets\emptyset$; $P\gets0$; $d\gets0$; $t_c\gets1$
\For{each frame $t=1,2,\dots$}
  \If{$t=1$ \textbf{or} $t \bmod 10=1$}
    \State $e_t,\beta_t,\nu_t\gets$ image cues of $I_t$
    \State $N,C,T\gets$ box cues of frame $t-1$
    \State add $C$ to $H$ (keep 7); $S\gets\mathrm{mean}(H)$
    \State $\ell\gets$ label (Eq.~\ref{eq:label}); $g\gets$ target (Eq.~\ref{eq:target})
    \State update $d$ and $P$ (Eqs.~\ref{eq:drop}--\ref{eq:peak})
    \If{$1\leq d\leq3$} $g\gets$ HIGH
    \EndIf
    \State apply the holds (Eq.~\ref{eq:hold})
    \If{$t-t_c\geq30$ \textbf{and} $g\neq q$}
      \State $q\gets g$; $t_c\gets t$
    \EndIf
  \EndIf
  \State $(r_t,\eta_t,c_t)\gets$ action of state $q$ (Table~\ref{tab:actions})
  \State detect with $(r_t,\eta_t)$; drop scores $<c_t$
  \State update the tracker; store its boxes
\EndFor
\end{algorithmic}
\end{algorithm}

\subsection{Action map}
Table~\ref{tab:actions} gives the operating point of each state. LOW uses the
largest input (1536 pixels), loose NMS (0.70), and a low score threshold
(0.10), so the detector passes more weak boxes when the scene is sparse.
MEDIUM uses 1280 pixels, tight NMS (0.45), and a score threshold of 0.40.
HIGH uses 1088 pixels, NMS 0.60, and a score threshold of 0.40. In the two
difficult states, the detector therefore passes only boxes with a score of at
least 0.40. The map was selected together with $\tau_m$ and $\tau_h$ by the
joint search on the calibration sequences (Sect.~\ref{sec:constants}). The
score threshold $c_t$ is applied on the detector side before the tracker. It is
separate from any confidence gate inside the tracker, such as the fixed
tracker-side gate of OATrack.

\begin{table}[t]
\caption{Action map. State names describe scene difficulty, not a resolution
rank. Selected by the joint search on the calibration sequences.}
\label{tab:actions}
\centering\footnotesize
\begin{tabular}{@{}lrrr@{}}
\toprule
State & Resolution $r$ & NMS IoU $\eta$ & Score threshold $c$\\
\midrule
LOW & 1536 & 0.70 & 0.10\\
MEDIUM & 1280 & 0.45 & 0.40\\
HIGH & 1088 & 0.60 & 0.40\\
\bottomrule
\end{tabular}
\end{table}

\subsection{How each value was chosen}\label{sec:constants}
All values of Tables~\ref{tab:actions} and~\ref{tab:constants} were fixed
before the final evaluation and were not changed after it. The last column of
Table~\ref{tab:constants} gives the source of each value. There are four
sources.

\emph{Joint search on the calibration sequences.} The SCI thresholds $\tau_m$
and $\tau_h$ and the action map were selected together by an Optuna TPE search
with 60 trials \cite{akiba2019optuna}. The calibration sequences are eight
sequences of the VisDrone2019-MOT training split. They serve as the validation
(tuning) data of the final stage and are disjoint from the evaluation
sequences. Each trial was scored by four-fold cross-validation, with two
sequences held out in each fold, as the mean change in HOTA against the
baseline. The search ranges were $\tau_m\in[0.10,0.50]$ and
$\tau_h\in[\tau_m+0.05,\min(\tau_m+0.40,0.90)]$, with, for each state, an
input of 1088, 1280, or 1536 pixels, an NMS value of 0.45, 0.60, or 0.70,
and a score threshold of 0.10, 0.25, or 0.40. The input sizes and NMS values
are the operating points for which detector outputs were cached; 1536 pixels
with NMS 0.70 is the baseline point, and 0.70 is the default NMS value of the
Ultralytics detector. The three score thresholds are the low and high
thresholds of ByteTrack (0.10 and 0.25) and the confidence gate of OATrack
(0.40). In a separate sweep of the score threshold over 0.01, 0.10, 0.25, and
0.40 on the same sequences, at 1088 pixels and NMS 0.45, 0.40 gave the highest
pooled HOTA (+0.881 over 0.01) and improved seven of the eight sequences. The frozen policy reached a mean HOTA
gain of 2.447 points over the four folds, with all folds positive
(Sect.~\ref{sec:modern}).

\emph{Cue audit on the calibration sequences.} The choice of the crowd cue as
the only SCI input came from testing all 31 non-empty subsets of the five
cues with the same cross-validation (Sect.~\ref{sec:modern}).

\emph{Temporal sweep on the validation split.} The analysis stride and the
SCI window were selected in the first development stage, with the YOLOv8n
pipeline, on the VisDrone2019-MOT validation sequences. The sweep tested 25
combinations: windows of 1, 3, 5, 7, and 9 analysis steps and strides of 1, 5,
10, 15, and 20 frames. Five of the 25 settings met both constraints of that
stage, at least 25 FPS and at most 271 IDS (the IDS of the original controller
on the same data). Among them, a window of 7 and a stride of 10 gave the
highest MOTA, 18.165 against 18.044 for the next setting. Operating-point
sweeps on the same data covered input resolutions from 512 to 960 pixels,
confidence thresholds from 0.05 to 0.50, and NMS thresholds from 0.30 to 0.80,
and defined the supported values of the first optimized profiles
(Sect.~\ref{sec:development}).

\emph{Fixed design values.} The tiny-object area follows the COCO definition
of small objects \cite{lin2014coco}. The remaining values (the crowd
normalizer, the Canny thresholds, the label thresholds, the recovery, hold,
and dwell values, and the initial state) come from the original AC-MOT
controller and were kept unchanged in every search of the final stage. The
searched values above
were tuned with these values in place, so the selected thresholds and actions
are matched to them. For example, the crowd normalizer of 30 only sets the
scale of $S_j$: $\tau_m$ and $\tau_h$ were searched on this scale and
correspond to about 9 and 15 tracked objects.

\begin{table*}[t]
\caption{Values of the AC-MOT controller and how each value was chosen.
Calibration sequences: eight VisDrone2019-MOT training sequences used as
tuning data, disjoint from the evaluation sequences. CV: four-fold
cross-validation.}
\label{tab:constants}
\centering\footnotesize
\setlength{\tabcolsep}{4pt}
\begin{tabular}{@{}l>{\raggedright\arraybackslash}p{50mm}l>{\raggedright\arraybackslash}p{62mm}@{}}
\toprule
Symbol & Meaning & Value & How the value was chosen\\
\midrule
$\tau_m$ & SCI threshold for MEDIUM, Eq.~\eqref{eq:target} & 0.29747709701135766 & joint TPE search on the calibration sequences, 60 trials, CV\\
$\tau_h$ & SCI threshold for HIGH, Eq.~\eqref{eq:target} & 0.5114928923560124 & same joint search\\
$(r,\eta,c)$ & operating point of each state & Table~\ref{tab:actions} & same joint search over $\{1088,1280,1536\}\times\{0.45,0.60,0.70\}\times\{0.10,0.25,0.40\}$\\
-- & SCI input cue & crowd only & cue audit of all 31 cue subsets on the calibration sequences, CV\\
-- & analysis stride & 10 frames & temporal sweep of 25 settings on the validation split; highest MOTA with $\geq$25 FPS and $\leq$271 IDS\\
$|H_j|$ & SCI window & 7 steps & same temporal sweep\\
-- & tiny-object area & $32\times32$ pixels & COCO small-object limit \cite{lin2014coco}\\
-- & crowd normalizer, Eq.~\eqref{eq:crowd} & 30 boxes & original controller; sets the scale on which $\tau_m$, $\tau_h$ were searched\\
-- & Canny thresholds (lower, upper) & 50, 120 & original controller; fixed in all final-stage searches\\
-- & night and blur labels, Eq.~\eqref{eq:label} & $\beta_t<80$, $\nu_t<180$ & original controller; fixed in all final-stage searches\\
-- & crowded label, Eq.~\eqref{eq:label} & $C_j>0.65$ or $e_t>0.13$ & original controller; fixed in all final-stage searches\\
-- & tiny label and HIGH guard, Eqs.~\eqref{eq:label}, \eqref{eq:target} & $T_j>0.50$ & original controller; fixed in all final-stage searches\\
-- & recovery: minimum peak, drop ratio, peak decay, steps & 5, 0.40, 0.80, 3 & original controller; fixed in all final-stage searches\\
-- & holds in HIGH ($S_j$, $T_j$) and MEDIUM ($S_j$), Eq.~\eqref{eq:hold} & (0.50, 0.40), 0.25 & original controller; fixed in all final-stage searches\\
-- & dwell, Eq.~\eqref{eq:dwell} & 30 frames & original controller; fixed in all final-stage searches\\
-- & initial state & LOW & original controller; fixed in all final-stage searches\\
\bottomrule
\end{tabular}
\end{table*}
